\documentclass{situationist-preprint}

\title{Simulated Physics and the Evidential Base}
\msid{MET-25-b}
\msfamily{MET-25}
\msdiscipline{Metaphysics}
\mstype{Companion manuscript}
\mskeywords{simulation argument; evidential base; fidelity; limits of computation; self-undermining arguments}

\begin{document}
\maketitle

\begin{abstract}
MET-25-a argues that the simulation argument is self-undermining because its feasibility premise is supported only by our physics, which is evidence about the base level of reality only if we are not simulated. The key premise is that, if we are simulated, our evidence is silent on the computing resources of the base level. We defend it in detail. We distinguish three relations a simulated physics may bear to the physics of its simulators, show that our evidence bears on base resources only through a fidelity premise, and argue that the argument's own reference class makes fidelity improbable given that we are simulated. Empirical signatures of simulation, if found, would support the hypothesis directly but would not rescue the argument.
\end{abstract}

\section{The evidential base of the feasibility premise}

The feasibility premise $P$ of MET-25-a (Definition~2.1) says that the physics of the base level permits a mature civilization to run very many ancestor-simulations. Our evidence for it is our physics: estimates of the computing power that a planet-sized computer could deliver, derived from physical limits on the speed and energy cost of computation \citep{Lloyd2000}, compared with estimates of the cost of simulating human brains and environments \citep{Bostrom2003}. These limits are laws of our physics. Assumption~2.4 of MET-25-a states that our evidence supports $P$ in a world only if we are not simulated in it. A defender of the argument may grant that the evidence is less direct if we are simulated and still deny that it is silent. We examine the possible relations between a simulated physics and the base physics.

\begin{definition}[Fidelity]\label{d:fid}
A simulation is \emph{faithful} if the laws governing its simulated world, including the physical limits on computation within it, are those of the world that runs it.
\end{definition}

Three cases are possible. In a faithful simulation, our physics is a copy of the base physics. In a coarse-grained simulation, the simulated world reproduces macroscopic regularities and simulates fine detail only where observers look, as Bostrom himself suggests efficient simulators would \citep{Bostrom2003}; its laws, as we find them, are partly an artefact of what is rendered. In an arbitrary simulation, the rules are chosen by the simulators and need resemble nothing in their own world. In the second and third cases, the physical limits on computation that we measure are limits of the simulation's rules, and they carry no information about base resources.

\begin{proposition}[Fidelity is required]\label{p:fid}
If we are simulated, our physics bears on $P$ only if we are in a faithful simulation.
\end{proposition}

\begin{proof}
If the simulation is coarse-grained or arbitrary, the laws we measure differ from the base laws in ways our evidence cannot reveal, so they fix no bound on what the base laws permit. Only if the simulated laws are the base laws does a bound we derive from ours hold of the base.
\end{proof}

\section{Fidelity given simulation}

The defender will say that ancestor-simulations are by definition faithful, since they simulate the simulators' own past. This is the first objection considered in MET-25-a, and we can now answer it in full. The simulation argument does not count only faithful simulations. Its indifference principle applies to all observers whose experiences are like ours, and its conclusion is that most such observers are simulated. A coarse-grained simulation of an ancestral world produces observers with experiences like ours at a fraction of the cost of a faithful one, and an arbitrary simulation that reproduces the appearances of our world does so too. By the argument's own reasoning, if simulators are cost-sensitive, most simulated observers like us are in simulations that are not faithful at the level of fundamental physics.

\begin{proposition}[Fidelity improbable]\label{p:improb}
If the simulations that produce observers like us are mostly coarse-grained or arbitrary, then, by the argument's indifference principle, the credence that we are in a faithful simulation, given that we are simulated, is low.
\end{proposition}

\begin{proof}
The indifference principle assigns equal credence to being any observer with experiences like ours. If most simulated observers with such experiences are in unfaithful simulations, the conditional credence in fidelity given simulation is the fraction of them in faithful simulations, which is small.
\end{proof}

Propositions~\ref{p:fid} and~\ref{p:improb} together support the silence assumption of MET-25-a (Assumption~3.5): given that we are simulated, our evidence is close to silent on $P$, since it bears on $P$ only through fidelity, and fidelity given simulation is improbable. The quantitative caps of MET-25-a then apply: by Proposition~3.6 and Corollary~3.7 of that manuscript, an argument-based credence in the simulation hypothesis cannot be raised above a modest level. \citet{Birch2013} reaches a related conclusion by a different route, arguing that the simulation argument requires a selective scepticism that it cannot justify; the present point locates the selectivity in the fidelity premise.

\section{Empirical signatures}

\citet{BeaneDavoudiSavage2014} proposed that a simulation of our universe on a space-time lattice might leave detectable traces, such as anisotropies in the distribution of the highest-energy cosmic rays. Suppose such traces were found. They would be evidence that we are simulated, and our credence in the simulation hypothesis would rise. But the evidence would concern the simulation's discretization, a feature of its rules, and would say nothing about the computing resources of the base level beyond the fact that it can run at least one simulation, ours. It would therefore leave $P$, which concerns very many simulations, unsupported.

This is consistent with MET-25-a, whose Theorem~3.2 concerns argument-based attitudes only: it denies that the simulation argument can make us confident that we are simulated, not that any evidence could. Independent evidence of simulation would bypass the argument rather than rescue it, and the lesson of the argument's failure would stand: our physics can tell us what our world permits, and it can tell us what the base level permits only on the assumption that we are at it.

\begin{thebibliography}{99}\small
\bibitem[Beane et~al.(2014)]{BeaneDavoudiSavage2014} Beane, S.~R., Davoudi, Z. and Savage, M.~J. (2014). Constraints on the universe as a numerical simulation. \emph{The European Physical Journal A}, 50, 148.
\bibitem[Birch(2013)]{Birch2013} Birch, J. (2013). On the \enquote{simulation argument} and selective scepticism. \emph{Erkenntnis}, 78, 95--107.
\bibitem[Bostrom(2003)]{Bostrom2003} Bostrom, N. (2003). Are you living in a computer simulation? \emph{Philosophical Quarterly}, 53, 243--255.
\bibitem[Lloyd(2000)]{Lloyd2000} Lloyd, S. (2000). Ultimate physical limits to computation. \emph{Nature}, 406, 1047--1054.
\end{thebibliography}

\end{document}
